\documentclass[10pt,a4paper]{amsart}
\usepackage[margin=19mm]{geometry}
\usepackage{amsmath,amssymb,mathtools}
\usepackage[scr=rsfs,cal=euler]{mathalfa}
\usepackage{xfrac}
\usepackage[unicode,hidelinks,bookmarks=false]{hyperref}
\newcommand{\ie}{i.e.\ }
\newcommand{\cf}{cf.\ }
\newcommand{\resp}{respectively\ }
\newcommand{\Subsection}{\subsection}
\providecommand{\nobreakdash}{\nobreak\hbox{-}}
\setlength{\emergencystretch}{2em}
\multlinegap0pt
\usepackage{bm}
\usepackage{bbm}
\usepackage{dsfont}
\usepackage{xspace}
\usepackage{cancel}
\usepackage{mathtools}
\usepackage{amsthm}
\makeatletter
\@ifundefined{thm}{\newtheorem{thm}{Thm}}{}
\@ifundefined{lemma}{\newtheorem{lemma}[thm]{Lemma}}{}
\makeatother
\def\er{\mathbb R}
\title[Diffeomorphic approximation of piecewise affine homeomorphis]{Diffeomorphic approximation of piecewise affine homeomorphisms}
\author{Daniel Campbell, Luigi D'Onofrio, Tom\'a\v{s} V\'itek}
\date{Source: arXiv:2507.02854v2 (CC BY 4.0)}
\begin{document}
\maketitle

\noindent\emph{Source and licence.} Diffeomorphic approximation of piecewise affine homeomorphisms. Version arXiv:2507.02854v2 (CC BY 4.0), distributed under the Creative Commons Attribution 4.0 International licence, as recorded in the arXiv licence metadata for this exact version and shown on the abstract page https://arxiv.org/abs/2507.02854v2. Full rights notice: licenses/arxiv-2507.02854-CC-BY-4.0.txt.\par\smallskip

\noindent\textit{Editorial scope.} Counted unit: the Lemma whose source label is \texttt{FaceLem} in the article (source0.tex lines 343--367, source label \texttt{FaceLem}), delivered together with the article's own complete original proof of it (source0.tex lines 368--447, one \texttt{proof} environment of 80 lines). No mathematics is written, repaired or re-derived: statement and proof are the article's own text line for line. Required context retained uncounted: Book (lines 248--250). Every definition, display and label of those lines is retained unchanged. Omitted from this package and not redistributed: 1--247, 251--342, 448--1030. Nothing inside a retained excerpt is omitted or altered: every retained body is delivered exactly as the article's lines are written, so no omission receipt of any kind accompanies this package. Citations: the retained excerpts carry no citation command at all, so no bibliography entry of the article is transcribed and no reference is redistributed. Reference adaptation: none was needed and none was made; every reference inside the delivered unit resolves inside it, so no unresolved reference or citation is produced. Printed slips kept literal: no printed word, symbol, hypothesis or number of the counted statement or of its proof was altered. Layout and class: the article's own class is \texttt{amsart} with packages including enumerate, graphicx, amsmath, verbatim, amssymb, pgf, tikz, inputenc, hyperref; the wrapper is the lane's pinned \texttt{amsart} editorial wrapper at 10pt on A4, which restates the theorem-like environments the retained text needs and adds the article's own macro definitions that the retained text needs (\texttt{\textbackslash er}), with extra wrapper packages (bm, bbm, dsfont, xspace, cancel, mathtools). The delivered numbering is the unit's own. Assets: no retained span contains a figure environment, a graphics-inclusion command, a PDF-page inclusion, a table or a TikZ picture; no image file is copied into this package, the package declares no asset dependency and no image file is redistributed with it. Licence: version arXiv:2507.02854v2 of this article is distributed under the Creative Commons Attribution 4.0 International licence, as recorded in the arXiv licence metadata for this exact version and shown on the abstract page https://arxiv.org/abs/2507.02854v2. A full rights notice is delivered at licenses/arxiv-2507.02854-CC-BY-4.0.txt. Characters: every retained excerpt is pure ASCII, so the delivered engine, XeLaTeX with its default Latin Modern text font, reports no missing character.
	\begin{lemma}\label{Book}
		Let $d\geq 1$, let $\Omega, \Delta \subset \mathbb{R}^d$ be bounded domains and let $\mu:\overline{\Omega}\to \mathbb{R}^d$ be a local diffeomorphism of $\Omega\subset \mathbb{R}^d$ and continuous on $\overline{\Omega}$. Further we assume that there is a neighborhood $O$ of $\partial\Omega$ in $\overline{\Omega}$ and a neighborhood $D$ of $\partial\Delta$ in $\overline{\Delta}$ such that $\mu$ is a homeomorphism $O$ onto $D$. Then $\mu$ is a diffeomorphism of $\Omega$ onto $\Delta$. 
	\end{lemma}

	\begin{lemma}\label{FaceLem}
		Let the affine mappings $A_1, A_2 : \mathbb{R}^3 \to \mathbb{R}^3$ be such that the mapping
		$$
		f(x)=\begin{cases}
			A_1(x) \quad &\text{for }x_1\leq0\\
			A_2(x) \quad &\text{for }x_1\geq 0.
		\end{cases}
		$$
		is a sense-preserving homeomorphism. Further, let $[\partial_1 A_2]^1 \geq [\partial_1 A_1]^1>0$ and $w\in \mathcal{C}^{\infty}(\er^2)$ be a bounded positive function. Then, there exists a $\sigma>0$ such that whenever $|Dw(y,z)|\leq \sigma$ the map 
		\begin{equation}\label{defg}
			g_w(x) = \Big(1-\eta\Big(\frac{x_1}{w(x_2,x_3)}\Big)\Big)A_1(x) + \eta\Big(\frac{x_1}{w(x_2,x_3)}\Big)A_2(x)
		\end{equation}
		is a diffeomorphism $g_w: \mathbb{R}^3 \to \mathbb{R}^3$ satisfying 
		\begin{equation}\label{NoChange}
			g_w(x_1,x_2,x_3) = f(x_1,x_2,x_3)
		\end{equation}
		on the set
		$$
		\{(x_1,x_2,x_3)\in \er^3; x_1\leq 0\}\cup \{(x_1,x_2,x_3)\in \er^3; x_1\geq w(x_2,x_3)\}.
		$$
		Further, there exists an absolute constant $C_{\ref{Guessed}}$ independent of $f$ and $w$ such that 
		\begin{equation}\label{Guessed}
			\|Dg_w\|_{L^\infty}\leq C_{\ref{Guessed}}(1+\sigma)\|Df\|_{L^\infty}.
		\end{equation}
	\end{lemma}

	\begin{proof}
		Since $A_1$, $A_2$, $w$ and $\eta$ are all in $\mathcal{C}^{\infty}$ and $w>0$ everywhere we have that $g_w\in \mathcal{C}^{\infty}(\er^3)$. Further, from the definition, it is immediately obvious that $g_w(x_1,x_2,x_3) = f(x_1,x_2,x_3)$ on the set $\{x_1\leq 0\}\cup \{x_1\geq w(x_2,x_3)\}$.
		
		
		Note that, since $A_1 = A_2$ on the plane $\{x\in \er^3; x_1=0\}$, we have $\partial_2A_1 = \partial_2A_2$ and $\partial_3A_1 = \partial_3A_2$. The definition of $g_w$ is independent of addition by a vector, i.e., let $v\in \er^3$ and let $\tilde{f} = f+v$, then $\tilde{g}_{w} =g_w + v$. The same holds regarding the action of linear isometries, $O(3)$. Therefore, without loss of generality, we can assume that $f$ maps the plane $\{x\in \er^3: x_1=0\}$ onto itself, while still having positive Jacobian and having $[\partial_1 A_2]^1 \geq [\partial_1 A_1]^1>0$. Our choice of isometry implies that $[\partial_2 A_1]^1 = [\partial_3 A_1]^1 = [\partial_2 A_2]^1 = [\partial_3 A_2]^1 = 0$.
		
		
		Since we have that
		\begin{equation}\label{TheMatrix}
			Dg_w = \left( \begin{matrix}
			[\partial_1g_w]^1 &0 &0\\
			[\partial_1g_w]^2 &[\partial_2g_w]^2 & [\partial_3g_w]^2\\
			[\partial_1g_w]^3 &[\partial_2g_w]^3 & [\partial_3g_w]^3
		\end{matrix}\right).
		\end{equation}
		we calculate the Jacobian of $g_w$ by an expansion of the first row. We have
		\begin{equation}\label{D1Posi}
			\begin{aligned}[]
				[\partial_1g_w]^1  =&  \Big(1-\eta\Big(\frac{x_1}{w(x_2,x_3)}\Big)\Big)[\partial_1 A_1]^1+ \eta\Big(\frac{x_1}{w(x_2,x_3)}\Big)[\partial_1 A_2]^1\\
				&  \quad +\frac{1}{w(x_2,x_3)}\eta'\Big(\frac{x_1}{w(x_2,x_3)}\Big)\big([A_2(x)]^1 - [A_1(x)]^1\big) \\
				\geq&  [\partial_1 A_1]^1 + \frac{1}{w(x_2,x_3)}\eta'\Big(\frac{x_1}{w(x_2,x_3)}\Big)\big([A_2(x)]^1 - [A_1(x)]^1\big) \\
				\geq&  [\partial_1 A_1]^1 >0.
			\end{aligned}
		\end{equation}
		Now we want to show two things. Firstly, that 
		$$
		\det \left( \begin{matrix}
			[\partial_2g_w]^2 & [\partial_3g_w]^2\\
			[\partial_2g_w]^3 & [\partial_3g_w]^3
		\end{matrix}\right)>0
		$$
		implying that $J_{g_w}>0$ everywhere. Secondly we want to prove that $g_w$ is equal to a homeomorphism outside of some large ball. Then we can finish the proof by applying Lemma~\ref{Book}. 
		
		
		Regarding the first point, note that at every $x$, we have $\partial_j A_1(x) = \partial_j A_2(x) = \partial_j f(x)$ for $ j = 2, 3$. Further, it holds that
		$$
			\begin{aligned}
				|\partial_jg_w(x) - \partial_jA_1| &\leq \partial_jw(x_2,x_3) \frac{x_1|A_2(x) - A_1(x)|}{w^2(x_2,x_3)}\eta'\Big(\frac{x_1}{w(x_2,x_3)}\Big).
			\end{aligned}
			$$
		Either $\eta'\Big(\frac{x_1}{w(x_2,x_3)}\Big) = 0$, or $x_1<w(x_2,x_3)$ implying that $|A_2(x) - A_1(x)|< 2\|Df\|_{\infty}w(x_2,x_3)$. Therefore, we have
		\begin{equation}\label{EstimateDJ}
				\begin{aligned}
				|\partial_jg_w(x) - \partial_jA_1|\leq 2\|Df\|_{\infty}|Dw|.
			\end{aligned}
		\end{equation}
		The above implies that
		\begin{equation}\label{TruEstimateDJ}
				\|\partial_jg_w(x)\|_{\infty}\leq C\|Df\|_{\infty}(1+\|Dw\|_{\infty})
		\end{equation}
		for an absolute constant $C$. By the continuity of the determinant on entries, it follows, for a sufficiently small choice of $\sigma$ and $|Dw|\leq \sigma$, that 
		\begin{equation}\label{Poker}
			\det \left( \begin{matrix}
			[\partial_2g_w]^2 & [\partial_3g_w]^2\\
			[\partial_2g_w]^3 & [\partial_3g_w]^3
		\end{matrix}\right) \geq \tfrac{1}{2}\det \left( \begin{matrix}
		[\partial_2f]^2 & [\partial_3f]^2\\
		[\partial_2f]^3 & [\partial_3f]^3
		\end{matrix}\right)>0.
		\end{equation}
		Thus we have concluded that $J_{g_w}>0$ everywhere.
		
		Instead of showing that $g_{w}$ is injective outside a large ball we define a modified $\tilde{w}$ equal to $w$ close to the origin  and show that the modified is injective outside a large ball. Let $r>0$ be given and let
		$$
			R:=r+2+\frac{\sup\{|w|\}+1}{\sigma}.
		$$
		Then we can define a smooth $\tilde{w}$ such that $\tilde{w} = w$ on $B(0,r)\subset \er^2$, $w= 1$ outside $B(0,R)$ and $|D\tilde{w}|\leq \sigma$. To achieve this let us firstly define $\tilde{\tilde{w}} := w$ on $B(0,r+1)\subset \er^2$, $w= 1$ outside $B(0,R-1)$ and extend this map onto $\er^3$ by the McShame extension. We define $\tilde{w}$ as the standard mollification of $\tilde{w}$. Since $g_w = g_{\tilde{w}}$ on $\er\times B(0,r)$ it suffices to show that $g_{\tilde{w}}$ is a diffeomorphism. Since we already know that the smooth map $g_{\tilde{w}}$ has positive Jacobian everywhere and is therefore a local diffeomorphism, it suffices to show that $g_{\tilde{w}}$ is injective outside $Q(0,R)\subset \er^3$.
		
		Since $R > |w|$ we have on the set $\big(\er^3\setminus Q_{\er^3}(0,R)\big) \cap \{x_1>1\}$ that $g_{\tilde w} = f$ and $f$ is injective. On the set $\er\times(\er^2\setminus Q_{\er^2}(0,R))$ we have, by \eqref{TheMatrix} and  \eqref{D1Posi}, that $g_{\tilde{w}}^1(x) = g_{\tilde{w}}^1(y)$ if and only if $x_1 = y_1$. Now, thanks to the fact that $\tilde{w}(x_2,x_3) =1$ for all $(x_2,x_3)\in \er^2\setminus Q_{\er^2}(0,R)$ we have
		\begin{equation}\label{Sunlight}
			\partial_2 g_w = \partial_2f \qquad \text{and}\qquad \partial_3 g_w = \partial_3f.
		\end{equation}
		The above and the injectivity of $f$ implies that the restriction of $g_{\tilde{w}}$ to hyperplanes of type $\{x\in \er^3; x_1 = c\}$ is injective. The combination of these facts imply that $g_{\tilde{w}}$ is injective on the set $\er\times(\er^2\setminus Q_{\er^2}(0,R))$. In total we see that $g_{\tilde{w}}$ is injective on $\er^3\setminus Q_{\er^3}(0,R)$. It follows from Lemma~\ref{Book} that $g_{\tilde{w}}$ is injective on $\er^3$, especially it is injective on $B(0,r)$ and so $g_{w}$ is a diffeomorphism of $B(0,r)$. On the other hand, $r>0$ was chosen arbitrarily and therefore $g_w$ is a diffeomorphism on $\er^3$. 
		
		We estimate $\partial_1g_w$ as follows;
		$$
		|\partial_1g_w| \leq  |\partial_1 A_1|+ |\partial_1 A_2| +\eta'\Big(\frac{x_1}{w(x_2,x_3)}\Big)\frac{|A_2(x) - A_1(x)|}{w(x_2,x_3)}.
		$$
		As pointed out before \eqref{EstimateDJ}, we either have $\eta' = 0$ or $|A_2 (x) - A_1(x)|\leq 2\|Df\|_{\infty}w(x_2,x_3)$. Thus we see $|\partial_1g_w|$ is bounded independently of the values of $w$. To conclude \eqref{Guessed} it suffices to consider \eqref{TruEstimateDJ}. 
	\end{proof}

\end{document}
